\chapter{相关理论与技术}

\section{基本概念}

介绍论文中涉及的基本概念和定义。本节阐述研究所需的理论基础，为后续研究奠定理论支撑。

\section{相关技术}

详细介绍论文中使用的相关技术。对每种技术的原理、特点和应用场景进行系统性阐述。

\subsection{技术A}

技术A的详细描述。包括该技术的基本原理、发展历程、主要特点以及在本研究中的应用价值。

\subsection{技术B}

技术B的详细描述。分析该技术与技术A的关系，以及两者在本研究中如何协同工作。

\section{示例：定理、证明与智能引用}

\begin{definition}\label{def:feature-space}
设样本集合为 $\mathcal{X}$，映射 $\phi: \mathcal{X}\rightarrow\mathbb{R}^{d}$ 所诱导的向量空间称为特征空间。
\end{definition}

\begin{theorem}\label{thm:energy-balance}
若系统质量 $m>0$ 且光速 $c$ 为常数，则系统静能满足
\begin{equation}
    E = mc^2.
    \label{eq:einstein}
\end{equation}
\end{theorem}

\begin{proof}
该式是质能等价关系的标准表达。模板在这里重点演示定理、证明和公式编号，而非物理推导本身。结合\cref{thm:energy-balance,eq:einstein}可以看到定理环境与智能交叉引用已经可直接使用。
\end{proof}

\section{示例：伪代码}

\begin{algorithm}[H]
    \caption{模型训练流程}
    \label{alg:training}
    \begin{algorithmic}[1]
        \Require 训练集 $\mathcal{D}$，初始参数 $\bm{\theta}_{0}$，学习率 $\eta$
        \Ensure 收敛参数 $\bm{\theta}^{*}$
        \State 初始化轮次 $t \gets 0$
        \While{未满足停止条件}
            \State 从 $\mathcal{D}$ 中采样一个小批量 $\mathcal{B}$
            \State 计算损失 $\mathcal{L}(\bm{\theta}_{t}; \mathcal{B})$
            \State 更新参数 $\bm{\theta}_{t+1} \gets \bm{\theta}_{t} - \eta \nabla \mathcal{L}(\bm{\theta}_{t}; \mathcal{B})$
            \State $t \gets t + 1$
        \EndWhile
        \State \Return $\bm{\theta}_{t}$
    \end{algorithmic}
\end{algorithm}

\section{示例：插入图片}

图片应放在\texttt{figures/}目录下。图片和表格中的文字使用五号楷体。示例：

\begin{figure}[H]
    \centering
    % \includegraphics[width=0.6\textwidth]{figures/example.png}
    \caption{示例图片标题}
    \label{fig:example}
\end{figure}

可以使用\figref{fig:example}来引用图片（注意：正文中引用时不需要"\texttt{}"符号）。

\section{示例：插入表格}

表格使用三线表格式，标题位于表格上方。

\begin{table}[H]
    \centering
    \caption{示例表格标题}
    \label{tab:example}
    \begin{threeparttable}
        \begin{tabular}{lSS}
            \toprule
            方法 & {推理时延/\si{\milli\second}} & {准确率/\si{\percent}} \\
            \midrule
            基线模型 & 23.5 & 85.3 \\
            改进模型 & 18.2 & 91.2 \\
            \bottomrule
        \end{tabular}
        \begin{tablenotes}[flushleft]
            \item 注：数值列采用 \texttt{siunitx} 对齐，表注由 \texttt{threeparttable} 提供。
        \end{tablenotes}
    \end{threeparttable}
\end{table}

可以使用\tabref{tab:example}来引用表格。

\section{示例：数学公式}

行内公式示例：$E=mc^2$，公式与正文在同一行内。

行间公式需要单独占一行并居中显示：
\begin{equation}
    \int_{a}^{b} f(x) \, dx = F(b) - F(a)
    \label{eq:example}
\end{equation}

可以使用\eqnref{eq:example}来引用公式。注意公式编号格式为“(章节号-序号)”。

\section{示例：代码环境}

\begin{lstlisting}[language=Python,caption={训练脚本片段},label={lst:trainer}]
def train_one_epoch(model, dataloader, optimizer):
    model.train()
    for batch in dataloader:
        loss = model(batch)
        loss.backward()
        optimizer.step()
        optimizer.zero_grad()
\end{lstlisting}

模板默认提供统一的代码样式，也可结合\cref{alg:training,lst:trainer}在正文中直接引用算法和代码。

\section{本章小结}

本章对××××进行了系统性阐述，并演示了定理、伪代码、代码、图表、公式、智能引用和单位排版等常见学术写作能力。
